\section{Inputs and normalisation}
\label{sec:inputs}

\subsection{Input structure}

Inputs are collected in six groups: project information; site characteristics;
climate and heat exposure; vulnerability and equity; the proposed NbS
intervention; and cost and energy parameters. The full field reference, with
required or optional status and the normalisation applied to each field, is
reproduced in \Cref{app:fields}.

Two structural changes arrive at version \methver{}. The intervention field
\fieldname{nbs_type} is now a \emph{list}: an assessment may propose a package
of simultaneously available entries, scored under \Cref{sec:packages}. And the
site-characteristics group gains four questions that decide which interventions
are offered and feed no score at all (\Cref{sec:availability}).

Every input is either \emph{required}, in which case the assessment cannot
proceed without it, or \emph{optional}, in which case its absence reduces
confidence but never blocks the assessment. The interaction design anticipates
roughly forty-five inputs, of which it expects about twenty to be mandatory;
the implemented input schema requires exactly the four fields that a
documented formula cannot be evaluated without, and \Cref{app:fields} records
the status and normalisation of every field at version \methver{}. Reducing the total count
would not reduce uncertainty; it would move uncertainty from visible --- a
stated confidence level --- to invisible --- a silent default. The interface
consequently asks everything, makes skipping free, and shows the user what
skipping costs.

\subsection{Qualitative normalisation}
\label{sec:qualitative}

Most inputs accept a qualitative level, because most users at screening stage
do not have measured data. Levels map to a \numrange{0}{100} scale as shown in
\Cref{tab:levels}.

\begin{table}[H]
\centering
\caption[Qualitative level normalisation]{Normalisation of qualitative input
levels. The standard scale applies where a higher level indicates higher risk
or higher quantity; the inverted scale applies where a higher level indicates
\emph{lower} risk.}
\label{tab:levels}
\small
\begin{tabular}{@{}lrr@{}}
\toprule
\textbf{Level} & \textbf{Standard scale} & \textbf{Inverted scale} \\
\midrule
Low        & 25  & 100 \\
Medium     & 50  & 50  \\
High       & 75  & 25  \\
Very high  & 100 & --- \\
Unknown    & 50  & 50  \\
\bottomrule
\end{tabular}
\end{table}

Three fields use scales of their own, because their natural vocabularies are
not the four standard levels. Soil availability maps
\emph{none}\,$\rightarrow$\,25, \emph{limited}\,$\rightarrow$\,50,
\emph{moderate}\,$\rightarrow$\,75, \emph{high}\,$\rightarrow$\,100. Irrigation
availability maps \emph{none}\,$\rightarrow$\,25,
\emph{occasional}\,$\rightarrow$\,50, \emph{reliable}\,$\rightarrow$\,100.
Current shade level maps \emph{very low}\,$\rightarrow$\,25,
\emph{low}\,$\rightarrow$\,50, \emph{medium}\,$\rightarrow$\,75,
\emph{high}\,$\rightarrow$\,100. In each case \emph{unknown} maps to 50.

\subsection{The inverted mapping for cooling access}
\label{sec:inverted}

The principal application of the inverted scale is access to cooling refuge,
measured by two fields --- access to cooled \emph{indoor} space and access to a
cool \emph{outdoor} refuge --- where in each case \emph{high} access maps to
25, \emph{medium} to 50, and \emph{low} to 100. Both enter the Vulnerability
Score (\Cref{sec:vulnerability}) through a single \emph{cooling access
deficit} term: low access indicates higher vulnerability, so it must produce a
higher score.

The justification is \textcite{reid2009}. In that factor analysis of ten heat
vulnerability variables, air-conditioning prevalence emerged as one of four
factors explaining more than \SI{75}{\percent} of total variance --- and,
importantly, as an \emph{independent} dimension. That independence is what
justifies giving cooling access its own term in the vulnerability formula
instead of allowing it to be absorbed into a general deprivation indicator: a
factor that emerges separately alongside the social and environmental
vulnerability factor is not a restatement of it. That last step is this
methodology's inference from what a factor analysis means, not a claim the
source itself makes.

Two qualifications should be read alongside this. The study population was
North American, and access to mechanical cooling has a different distribution
and a different meaning in cities where air conditioning is rare across all
income groups. And the tool asks about access to cooled indoor space generally
rather than about air-conditioning ownership specifically, which is a
deliberate broadening --- a shaded, ventilated, publicly accessible cool refuge
counts --- but a broadening that the underlying study does not itself
validate.

The second field, added at version \texttt{2026.08.03}, addresses part of the
first qualification directly. Where mechanical cooling is rare, refuge from
heat is found outdoors, and the evidence for that is its own:
\textcite{burkart2016} and \textcite{sera2019} measure reachable vegetation and
water as modifiers of heat mortality rather than inferring them. The
combination rule for the two indicators, and the reason an unanswered indicator
is excluded rather than averaged in at the neutral 50, are given with the
formula in \Cref{sec:vulnerability}.

\subsection{The treatment of \emph{unknown}}
\label{sec:unknown}

Assigning \emph{unknown} a neutral score of 50 is a deliberate and disclosed
choice. Three alternatives were available and each carries a bias.

Treating an unknown as the best case flatters sites with poor data; treating it
as the worst case penalises them; and imputing a value from other fields
introduces a correlation the user never asserted and cannot see. The neutral
midpoint avoids both optimistic and pessimistic bias in the score, at the cost
of pulling every score toward the middle in proportion to how much is unknown.

That cost is not concealed. The confidence mechanism of \Cref{sec:confidence}
records, per output block, how much of the relevant input set was actually
supplied, so a score dominated by neutral defaults is reported at low
confidence. The two mechanisms are complementary: the score says what the
evidence supports, and the confidence says how much evidence there was. Neither
alone is sufficient, and quoting a score without its confidence rating is
identified in \Cref{sec:misuse} as a misuse of the tool.

\subsection{Land surface temperature as an optional proxy}
\label{sec:lst}

Where a user can supply a land-surface-temperature (LST) anomaly --- degrees
Celsius above the city mean, typically derived from satellite imagery --- the
tool converts it linearly to a \numrange{0}{100} score, with
\qty{0}{\degreeCelsius} mapping to 0 and \qty{+10}{\degreeCelsius} mapping to
100, clamped at both ends:

\begin{equation}
\label{eq:lst}
S_{\mathrm{LST}} \;=\; \clamp{10 \times a_{\mathrm{LST}}},
\end{equation}

where $a_{\mathrm{LST}}$ is the supplied anomaly in degrees Celsius and
$\clampop$ bounds its argument to $[0, 100]$. The \qty{10}{\degreeCelsius}
ceiling is consistent with observed intra-urban surface anomaly ranges while
keeping the normalisation interpretable.

\begin{keynote}
\noindent\textbf{This proxy is used under a stated restriction, and the
restriction is the most important sentence in this subsection.} Land surface
temperature is not air temperature, and the two diverge by a large and
systematic factor.

\textcite{venter2021} contrasted crowdsourced air temperature measurements
against satellite estimates and found the satellite-derived \emph{surface}
urban heat island to substantially exceed the air-temperature
(\emph{canopy-layer}) urban heat island: a reported mean surface UHI of
\qty{1.45}{\degreeCelsius} against a canopy-layer UHI of
\qty{0.26}{\degreeCelsius}, an approximately sixfold difference. This reflects
the physical distinction between surface, canopy-layer, and boundary-layer heat
islands established by \textcite{oke1976}.

The tool therefore treats LST \textbf{only as a relative indicator of which
parts of a city are hotter than others}. It is never interpreted as an
air-temperature value, never converted into a \si{\degreeCelsius} outcome, and
never compared against the tool's own air-temperature reduction estimates. It
influences only the ranking of sites within the Heat Exposure Score.
\end{keynote}

The verification status of \textcite{venter2021} is relevant here: the
bibliographic record was confirmed from the ADS record and an author-hosted
PDF, and the specific surface-versus-canopy comparison was confirmed from an
indexed secondary description, publisher full text being blocked. The claim the
methodology rests on --- that surface UHI substantially exceeds canopy-layer
UHI --- is corroborated conceptually by \textcite{oke1976}, and the tool's use
of the finding is qualitative (treat LST as relative only) rather than
numerical, so no tool output depends on the precise sixfold ratio.

\subsection{The four availability questions: inputs that gate but never score}
\label{sec:availability}

Version \methver{} adds four site-characteristic questions. They are unlike
every other input in this section, and the difference is the point of this
subsection.

\begin{table}[H]
\centering
\caption[The four availability questions]{The four availability questions added
at version \methver{}. Each decides which of the \num{121} catalogue entries
the tool offers for this site. None of them appears in any formula in
\Cref{sec:formulas}, in any confidence block of \Cref{sec:confidence}, or in
any aggregation.}
\label{tab:availability}
\small
\begin{tabularx}{\textwidth}{@{}l>{\raggedright\arraybackslash}X
                               >{\raggedright\arraybackslash}p{34mm}@{}}
\toprule
\textbf{Field} & \textbf{Answers} & \textbf{Gates} \\
\midrule
\fieldname{includes_railway} & yes / no &
Exactly one entry, the railway green corridor \\
\fieldname{existing_woodland} & yes / no &
Woodland \emph{restoration} entries only \\
\fieldname{waterfront_type} &
lake / river / dry river / covered river / sea &
Twelve water-landscape and ecological-network entries, by category \\
\fieldname{productive_governance} &
multi-select over community / individual / institutional / commercial &
The seven productive entries, by who delivers them \\
\bottomrule
\end{tabularx}
\end{table}

\begin{keynote}
\noindent\textbf{Gating is not scoring.} These four fields decide which
interventions the tool \emph{offers}, and they feed no score.

Making a gating answer a scoring input would need its own justification and
would move every existing result. The engine's test suite asserts the
separation as a property --- the same site scored with all four answered and
with none answered produces byte-identical output --- so the distinction cannot
erode quietly through a later edit. The fields are recorded in the input schema
and disclosed as availability-only in \Cref{tab:app-fields}.
\end{keynote}

\paragraph{What an unanswered question means differs by kind, deliberately.}
Three of the four describe a \emph{physical fact} about the site, and they gate
on positive confirmation: an entry that acts on an existing river is not
offered until the user says there is a river, because offering it would be
offering an intervention that cannot be built there. The fourth describes
\emph{who could deliver} the intervention, and it subtracts nothing until it is
answered.

That asymmetry is deliberate. The question originally specified for productive
landscapes was a yes/no --- is there interest in creating productive landscape
from individuals or communities --- which on a ``no'' would have suppressed the
urban farm and the agroforestry system: the highest-canopy, highest-cooling
entries in the group, and the ones a municipality is most likely to deliver. A
multi-select over four delivery models gates correctly at no extra cost in user
effort, and an \emph{empty} selection is not evidence that no delivery model
exists. The tool does not assert a negative from absent information, here as in
the soil--water pair of \Cref{sec:adjustment} and the suitability rules of
\Cref{tab:suitability-rules}.

Two consequences follow from the same principle and are worth stating because
both run against intuition. \textbf{Woodland \emph{creation} entries are
ungated}: urban woodland, microforest, afforestation, and the woodland buffer
are offered whatever the answer, because creating woodland does not require
woodland to be there already --- only restoration does. And \textbf{the four
constructed water features carry no waterfront condition at all}: a constructed
wetland, retention pond, detention basin, or water square needs no existing
water body, so gating them on one would withhold exactly the interventions
available to a site that has no water.

\paragraph{Land use maps to interventions rather than naming one.} The same
mechanism carries the retirement of \Cref{sec:retirements}. \num{77} of the
\num{121} entries name \emph{school} in their land-use list, \num{98} name
\emph{campus}, \num{81} \emph{healthcare}, and \num{44} \emph{memorial};
the assessment scale then narrows the menu further, so a school \emph{site}
is offered \num{71}. All four land-use totals and the nine published
situation counts are asserted by the test suite, so a land use silently losing
its interventions fails the build. Memorial is deliberately the sparsest: tree groves, meadows, woodland,
and quiet permeable planting are in; plazas, playgrounds, tram corridors,
living walls, and productive landscapes are out. All four counts are computed
with the governance question unanswered, which is the governance-agnostic
reading; answering it reduces them, and both readings are correct because they
answer different questions.

\subsection{Where the tool has no LST}

Most screening-stage users will not have an LST anomaly. The methodology
therefore provides a second path through the Heat Exposure Score
(\Cref{sec:heat-exposure}) that requires only a qualitative heat exposure level.
The two paths are not equivalent in information content, and the confidence
mechanism reflects that. The data-poor path is not a degraded mode to be
avoided; it is the expected mode for most assessments, and the tool is designed
to give a defensible answer in it.

\subsection{Inputs a map may derive, and the many it may not}
\label{sec:map-derived-inputs}

Version \methver{} adds an optional map step at the head of the questionnaire.
It is skippable in one action and the questionnaire is fully usable without
ever opening it; nothing in this subsection applies to an assessment that did
not use it.

Exactly three inputs can be derived from a point or a polygon on a map without
inventing anything. \texttt{site\_area\_m2} is pure geometry: the geodesic area
of the ring the user drew, evaluated on the authalic sphere, carrying no
uncertainty beyond the drawing itself. \texttt{country} is a point-in-polygon
test against public-domain boundaries \parencite{naturalearth}, deterministic given
the data, and feeds only the emission-factor and currency defaults. And
\texttt{climate\_zone} is a cited lookup plus a documented mapping: the
classification is published \parencite{beck2023}, and the mapping of its thirty
classes onto this tool's six zones is a methodology value in its own right,
derived below.

Everything else the questionnaire asks about the site is \emph{not} derived
from the map, and the reason is the whole point of this subsection. Canopy
cover, green cover, imperviousness, LST anomaly, land use, shade level,
population density and vulnerable-population presence all require satellite or
census data. The temptation is worth naming because it is strong: a map that
filled in canopy cover from a quick NDVI lookup would demonstrate well. It
would also produce this tool's \emph{most decision-relevant site inputs} from
an unvalidated pipeline, with no evidence table behind it and no confidence
treatment---which is precisely the defect the methodology refused when it
declined to ship default cost values (\Cref{sec:costs}), appearing
in a new place. The restriction is enforced at the service boundary rather than
left to the interface, so it holds as the questionnaire grows.

\subsubsection*{Mapping K\"oppen--Geiger onto six zones}

The classification is not ours; the mapping is. Thirty classes onto six zones
is a judgement that selects a row of the climate adjustment matrix
(\Cref{sec:adjustment}) and therefore changes results, so it is declared
as a methodology value with its own citation, derivation and version rather
than buried in the picker's code. It follows four principles, and every one of
the thirty assignments records which principle produced it, because a
thirty-row table of individual opinions could not be reviewed.

\begin{description}[leftmargin=1.5em, style=sameline, font=\normalfont\itshape]
  \item[Tropical] Group A splits on its dry season, which is the same split the
  tool's two tropical zones name: \texttt{Af} and \texttt{Am} to
  \texttt{tropical\_wet}, \texttt{Aw} to \texttt{tropical\_dry}.

  \item[Arid] Group B splits on desert (BW) versus steppe (BS) and the tool on
  arid versus semi-arid; these are the same axis. The hot/cold third letter is
  deliberately ignored, because the tool's arid downgrade encodes water
  limitation constraining evapotranspiration, not temperature---a cold desert
  limits evapotranspiration as a hot one does.

  \item[Warm season] Groups C and D whose third letter is \texttt{a} or
  \texttt{b} map to \texttt{temperate}. This is where the mid-latitude cooling
  literature was measured. Group~D is \emph{not} sent to \texttt{other}
  wholesale: Chicago is \texttt{Dfa}, Seoul \texttt{Dwa} and Moscow
  \texttt{Dfb}, all three have serious urban heat problems, and all three sit
  inside the evidence base behind the temperate row.

  \item[No warm season] Groups C and D whose third letter is \texttt{c} or
  \texttt{d}, and all of group E, map to \texttt{other}---fewer than four
  months above \SI{10}{\degreeCelsius}, or none. There is no warm season for
  the cooling literature to describe.
\end{description}

Mapping the polar and boreal classes to \texttt{other} rather than to
\texttt{temperate} matters more than it may appear. \texttt{other} is already
the tool's neutral climate condition---every family resolves to condition
\emph{unknown}, factor \num{1.0}---so an unclassifiable location receives
neither a boost nor a penalty. Forcing them into \texttt{temperate} would
assert that a tundra site cools like a temperate one, which no source in this
bibliography supports, and \emph{an uncited claim introduced through a lookup
table is still an uncited claim}.

The bundled classification layer is \ang{0.1}, roughly \SI{11}{\kilo\metre}.
That resolution is stated in the interface and in the report wherever a value
derived from it appears: the classification is reliable for the city a site
sits in, not for the site itself, which is the resolution at which climate zone
enters this methodology in any case, since it selects a row of an adjustment
matrix rather than a temperature.

\subsubsection*{An autofilled value counts as supplied, and says so}

The confidence mechanism (\Cref{sec:confidence}) counts what the tool was
\emph{told}, not how hard the user worked to tell it. A K\"oppen--Geiger lookup
at a location the user chose is a cited classification, not a guess, so it
counts as an answer. Not counting it would mean that clicking the map and
typing the same climate zone by hand produced \emph{different confidence
readings for identical inputs}, teaching users that the confidence meter
measures effort rather than information.

The value therefore counts and the provenance is disclosed instead: beside the
field in the interface, in the stored input, and itemised in the report's
Inputs sheet, which distinguishes user-supplied, autofilled and defaulted
values. Autofill never overwrites an answer the user has already given, and
overriding a value removes its marking. In practice none of the three fields
appears in any confidence block, so the question is one of principle rather
than arithmetic---no confidence figure moves either way.
